% chapters/chapter1.tex
\section{АНАЛІЗ ПРЕДМЕТНОЇ ОБЛАСТІ ТА ПОСТАНОВКА ЗАДАЧІ}

% Розділ 1 згідно нових методичних вказівок (10-15 стор.):
% 1.1 Характеристика предметної області
% 1.2 Огляд сучасних підходів та рішень
% 1.3 Системний аналіз задачі
% 1.4 Етичні, правові та регуляторні аспекти
% 1.5 Постановка задачі та формулювання вимог
% 1.6 Висновки до розділу 1
%
% Математичне та алгоритмічне обґрунтування (ДПФ, radix-4, DFG-модель,
% матричні факторизації) перенесено до розділу 2 «Проєктування».

% ─────────────────────────────────────────────────
\subsection{Характеристика предметної області}

\subsubsection*{Контекст застосування}

Швидке перетворення Фур'є (ШПФ, англ.\ \textit{Fast Fourier Transform}
--- FFT) є одним із найбільш вживаних обчислювальних алгоритмів
цифрової обробки сигналів (ЦОС). Воно лежить в основі цифрової
телекомунікації (OFDM-модуляція у стандартах LTE, 5G, Wi-Fi),
обробки мови й звуку, радіолокації, медичного томографування,
спектрального аналізу та сучасних алгоритмів машинного навчання
з перетворенням вхідних даних у частотний простір~\cite{nussbaumer1982,protsko2025}.
За обсягом щоденно виконуваних операцій ШПФ залишається однією з
найзавантаженіших обчислювальних процедур у світі.

Пряме обчислення дискретного перетворення Фур'є (ДПФ) має
квадратичну складність $O(N^2)$. Навіть для скромного розміру
$N = 1024$ це означає понад мільйон комплексних множень,
що неприйнятно для задач реального часу. Клас швидких алгоритмів,
започаткований Дж.~Кулі та Дж.~Тюкі у 1965~р.~\cite{cooley1965},
знижує складність до $O(N \log_2 N)$ завдяки рекурсивній
факторизації базисної матриці~\cite{nussbaumer1982}. Для довжин,
що є степенями чотирьох ($N = 4^t$), найефективнішим з точки зору
кількості нетривіальних множень є алгоритм з основою~4 (radix-4),
який шляхом двовимірної декомпозиції зводить одновимірне ДПФ
довжиною $N = 16$ до тристадійної схеми $4 \times 4$: чотири
паралельних ДПФ-4, множення на матрицю поворотних множників
$W_{16}^{rk}$ та знову чотири паралельних ДПФ-4~\cite{burrus1985dft}.

Паралельна природа алгоритму природно вкладається у парадигму
обчислень під керуванням потоку даних (англ.\ \textit{dataflow
computing}), теоретичні засади якої були закладені Дж.~Деннісом
(Jack B.\ Dennis) у MIT на початку 1970-х~рр.~\cite{dennis1974,dennis1975}.
У dataflow-моделі відсутній лічильник команд фон-нейманівської
архітектури: вузол графа активується тоді і лише тоді, коли на
всіх його вхідних каналах є токени, після чого самопланування
системи (англ.\ \textit{self-scheduling}) забезпечує максимальний
паралелізм без явних механізмів синхронізації~\cite{veen1986}.
Граф алгоритму ШПФ, що формально є орієнтованим ациклічним графом
залежностей за даними, відображається на граф потоку даних
(англ.\ \textit{Data Flow Graph} --- DFG) без перетворення моделі.

\subsubsection*{Опис проблеми та її актуальність}

Незважаючи на фундаментальну роль ШПФ та широкий спектр його
виробничих реалізацій, у навчальній та дослідницькій практиці
існує системна прогалина між трьома рівнями опису алгоритму:

\begin{itemize}
  \item \textit{математична модель} --- матричні факторизації,
    формули метеликів і поворотних множників, що подаються у
    підручниках із ЦОС~\cite{nussbaumer1982,burrus1985dft};
  \item \textit{формальна потокова модель} --- DFG і Synchronous
    Dataflow (SDF) Лі та Месершмітта, що описують вузли, дуги,
    токени та правило активації~\cite{lee1994,veen1986};
  \item \textit{виробнича реалізація} --- бібліотеки FFTW,
    ARM~CMSIS-DSP, NVIDIA~CuFFT, SPIRAL, Intel~IPP, що оптимізовані
    виключно під максимальну продуктивність на конкретних
    платформах~\cite{frigo2005,cmsis_dsp,du2012,puschel2004,intel_ipp}.
\end{itemize}

Ці рівні існують відокремлено один від одного. Математичні формули
не мають виконуваного аналога з можливістю поточного спостереження;
формальна SDF-модель рідко реалізується як самостійний артефакт;
а виробничі бібліотеки не надають жодних засобів для інтерактивного
дослідження внутрішньої потокової природи алгоритму: процесу
породження та поглинання токенів, динаміки завантаження вузлів,
часу наскрізної затримки та паралельних шляхів у графі. Студент,
викладач або інженер-початківець, що вивчає архітектуру потокового
обчислювача ШПФ, не має інструменту, який одночасно є математично
достовірним (виконує реальний алгоритм), формально обґрунтованим
(відповідає SDF-моделі) і візуально інтерактивним (показує рух
токенів у реальному часі).

Актуальність розробки візуального симулятора потокового обчислювача
ШПФ зумовлена трьома чинниками: (1)~зростанням попиту на навчальні
засоби з інтерактивною візуалізацією обчислювальних процесів у
підготовці ІТ-фахівців; (2)~розширенням інтересу до dataflow- й
потокових архітектур у сучасних спеціалізованих прискорювачах
(зокрема, тензорних процесорах і FPGA-рішеннях для ЦОС);
(3)~відсутністю вітчизняних веб-орієнтованих інструментів із
відкритим кодом для цього класу задач.


% ─────────────────────────────────────────────────
\subsection{Огляд сучасних підходів та рішень}

Аналіз літературних джерел і наявних програмних рішень за останні
10~років показує, що простір сучасних засобів обчислення ШПФ
поділяється на три сегменти: (а)~програмні бібліотеки загального
призначення; (б)~апаратні та GPU-прискорювачі; (в)~системи
автоматичної генерації оптимізованого коду. Паралельно існує
окрема галузь методів візуального моделювання, які застосовуються
для формалізованого опису обчислювачів, однак не реалізуються як
виконувані симулятори. Нижче розглядаються обидва ці напрямки та
проводиться їх порівняльний аналіз.

\subsubsection*{Програмні бібліотеки обчислення ШПФ}

\textbf{Бібліотека FFTW (Fastest Fourier Transform in the West)}
розроблена М.~Фріго і С.~Джонсоном у MIT у 1997~р.~\cite{frigo1998}
і у~1999~р. відзначена премією Вілкінсона. Архітектура FFTW
ґрунтується на трьох компонентах~\cite{frigo2005}: \emph{коделети}
--- компактні блоки прямолінійного C-коду для ДПФ невеликих розмірів,
автоматично згенеровані компілятором \texttt{genfft}; \emph{планувальник}
--- runtime-компонент, що методом динамічного програмування добирає
оптимальну рекурсивну факторизацію під конкретну апаратну платформу;
\emph{екзекутор} --- виконує складений план через явну рекурсію
«розділяй і володарюй». Рекурсивна декомпозиція є кеш-незалежною
(cache-oblivious): щойно підзадача поміщається у L1/L2-кеш, вона
виконується без промахів кешу, що і забезпечує вищу продуктивність
проти ітеративних реалізацій~\cite{frigo2005}.

\textbf{ARM CMSIS-DSP} --- стандартизований пакет ЦОС для
мікроконтролерів ARM Cortex-M~\cite{cmsis_dsp}. Надає жорстко
оптимізовані реалізації radix-4 для розмірів $N = 16, 64, 256,
1024, 4096$ із ручною адаптацією під SIMD-інструкції NEON і
DSP-розширення Cortex-M4/M7. Бібліотека не переноситься на інші
архітектури, не підтримує адаптивного планування та не надає
будь-яких засобів моніторингу стану обчислення.

\textbf{NVIDIA CuFFT} та \textbf{OpenCL FFT} --- адаптації ШПФ
для графічних процесорів. Стандарт OpenCL забезпечує функціональну
портативність, але не гарантує продуктивності: ядро, оптимізоване
під архітектуру NVIDIA~Fermi, демонструє низьку ефективність на
GPU ATI~Radeon через відмінності в ієрархії пам'яті та механізмах
планування потоків (warps/wavefronts)~\cite{du2012}. Для подолання
цього застосовується автотюнінг: А.~Нукада та С.~Мацуока~\cite{nukada2009}
показали, що адаптивний вибір відступів у спільній пам'яті та
кількості активних блоків потоків дозволяє перевершити навіть
вендорну бібліотеку NVIDIA~CuFFT.

\textbf{SPIRAL} (Carnegie Mellon University) --- принципово інший
підхід: замість вибору між написаними ядрами, SPIRAL генерує
платформо-залежний код безпосередньо з математичних формул
факторизації, описаних мовою SPL (Signal Processing Language).
Пошук оптимального графа зводиться до задачі пошуку у просторі
еквівалентних матричних розкладів, що може вирішуватися евристично
або методами машинного навчання~\cite{puschel2004}. SPIRAL
генерує код як для CPU (SIMD, AVX), так і для FPGA.

\textbf{Intel IPP (Integrated Performance Primitives)} займає
проміжне місце: жорстко оптимізовані примітиви ШПФ під векторні
інструкції SSE/AVX/AVX-512 процесорів Intel без автоматичного
пошуку факторизації під час виконання~\cite{intel_ipp}.

\subsubsection*{Методи візуального моделювання обчислювальних систем}

Окрему групу джерел утворюють методи візуального моделювання,
які теоретично описують структуру обчислювачів, але не реалізуються
як виконувані симулятори з числовою верифікацією.

\textbf{Загальні нотації моделювання} --- UML~Activity,
IDEF0/IDEF3, DFD Йордана/ДеМарко --- орієнтовані на control-flow
парадигму~\cite{rumbaugh2004,yourdon1979}. Вони не мають формальної
семантики токенів і не забезпечують однозначного відображення
паралелізму за даними, тому для потокового обчислювача ШПФ
породжують надлишкові конструкції \texttt{fork}/\texttt{join}.

\textbf{DFG/SDF} (Data Flow Graph і Synchronous Dataflow) ---
спеціалізовані моделі, розроблені для задач ЦОС~\cite{veen1986,lee1994,mit_dataflow}.
У SDF-моделі кількість токенів, що споживаються та виробляються
кожним вузлом за одну активацію, є фіксованою і відомою на етапі
статичного аналізу. Це дозволяє виконувати статичне планування без
ризику взаємних блокувань, гарантувати детермінованість результатів
і перевіряти консистентність графа алгебраїчно через топологічну
матрицю~\cite{schaumont2019}. Алгоритм ШПФ із фіксованою довжиною
$N = 16$ є природним прикладом SDF-графа.

\textbf{Сигнальні графи (signal flowgraphs) і butterfly-діаграми}
--- С.~С.~Буррус формалізує зв'язок між матричною факторизацією та
графом алгоритму через рівняння $\mathbf{X'} = \mathbf{C} \cdot
\mathbf{D} \cdot \mathbf{A} \cdot \mathbf{X}$, де кожна матриця
відповідає окремому ярусу графа~\cite{burrus1985dft}. Ці діаграми
є статичними і не моделюють динаміку виконання. У вітчизняній
науковій традиції аналогічний підхід застосовувався у роботах
ОКБ~ЕІВТ (М.~Розенблат, 1970-ті~рр.) та у монографії М.~М.~Яцимірського
(1997~р.)~\cite{protsko2025}. Зокрема, у Фізико-механічному
інституті ім.~Г.~В.~Карпенка НАН~України в 1980-х роках було
розроблено векторний спецпроцесор ШПФ на основі саме dataflow-парадигми
з блоком арифметики, що апаратно реалізовував чотири базові
потокові операції~\cite{protsko2025}.

\subsubsection*{Ступінь дослідженості проблеми та вузькі місця}

Узагальнення розглянутих джерел дозволяє зафіксувати такі
прогалини. По-перше, виробничі бібліотеки (FFTW, CMSIS-DSP,
CuFFT, SPIRAL, Intel~IPP) оптимізовані під продуктивність і не
надають жодних інтроспективних засобів: FFTW розкриває план
факторизації лише у текстовому форматі, жодна з бібліотек не
генерує анімований DFG і не показує рух токенів. По-друге, теоретичні
моделі DFG/SDF формально визначені, проте рідко супроводжуються
веб-доступним симулятором із числовою верифікацією. По-третє,
статичні засоби (draw.io, Visio, Simulink) або не відтворюють
динаміки, або потребують комерційних ліцензій, що обмежує
відтворюваність у навчальному та дослідницькому контекстах.

Таким чином, більшість авторів вирішує задачу обчислення ШПФ або
суто продуктивнісним підходом (FFTW, CuFFT, SPIRAL), або суто
теоретичним (DFG/SDF-моделі, сигнальні графи), однак підхід, який
поєднує формальну SDF-модель, виконуваний алгоритм із числовою
верифікацією та інтерактивну веб-візуалізацію, у літературі не
представлений.

\subsubsection*{Порівняння існуючих рішень}

Порівняльна характеристика розглянутих рішень за ключовими
характеристиками наведена у табл.~\ref{tab:solutions}.

\begin{table}[H]
  \caption{Порівняння існуючих рішень для обчислення ШПФ}
  \label{tab:solutions}
  \centering
  \small
  \begin{tabular}{p{2.4cm}p{2.0cm}p{1.8cm}p{1.6cm}p{1.8cm}p{2.4cm}}
    \toprule
    Рішення & Тип & Платформа & Динамічна візуалізація & Відкритий доступ & Ключова особливість \\
    \midrule
    FFTW 3.x & SW бібліотека & CPU (x86, ARM) & Немає & Так & Адаптивний planner, cache-oblivious рекурсія~\cite{frigo2005} \\[2pt]
    ARM CMSIS-DSP & SW бібліотека & ARM Cortex-M & Немає & Так & Фіксований radix-4, SIMD-оптимізація~\cite{cmsis_dsp} \\[2pt]
    NVIDIA CuFFT & SW бібліотека & GPU (NVIDIA) & Немає & Ні & Масовий паралелізм, вендор-лок~\cite{du2012} \\[2pt]
    Nukada autotuner & SW + автотюнінг & GPU (CUDA) & Немає & Част. & Runtime-компіляція~\cite{nukada2009} \\[2pt]
    SPIRAL & Генератор коду & CPU / FPGA & Немає & Так & Пошук у просторі факторизацій SPL~\cite{puschel2004} \\[2pt]
    Intel IPP & SW бібліотека & Intel x86 & Немає & Ні & Примітиви AVX/SIMD~\cite{intel_ipp} \\[2pt]
    MATLAB Simulink (SDF) & Середовище моделювання & Desktop & Часткова & Ні & SDF-модель, комерційна ліцензія \\[2pt]
    draw.io / Visio & Діаграмний редактор & Desktop / Web & Ні (статика) & Част. & Лише статична схема DFG \\[2pt]
    \textbf{Симулятор (ця робота)} & Веб-застосунок & Браузер & \textbf{Так} & \textbf{Так} & \textbf{Інтерактивний DFG + числова верифікація} \\
    \bottomrule
  \end{tabular}
\end{table}

Як видно з таблиці, жодне з розглянутих рішень не поєднує одночасно
достовірного числового обчислення, інтерактивної візуалізації руху
токенів та веб-доступності без комерційної ліцензії. Саме ця
незаповнена ніша визначає предмет цієї роботи та обумовлює
формулювання задачі у підрозділі~1.5.


% ─────────────────────────────────────────────────
\subsection{Системний аналіз задачі}

Результати підрозд.~1.1--1.2 підтверджують наявність об'єктивного
розриву між математичною моделлю, формальним DFG/SDF-описом та
виробничими реалізаціями алгоритму ШПФ. Для структурованого опису
цього розриву та переходу до постановки задачі застосовується
системний аналіз: визначаються зацікавлені сторони, цілі розробки,
обмеження та критерії успіху.

\subsubsection*{Зацікавлені сторони (stakeholders)}

Потенційні користувачі та бенефіціари розроблюваного симулятора
згруповані у табл.~\ref{tab:stakeholders}.

\begin{table}[H]
  \caption{Зацікавлені сторони та їхні основні потреби}
  \label{tab:stakeholders}
  \centering
  \small
  \begin{tabular}{p{3.4cm}p{5.0cm}p{5.0cm}}
    \toprule
    Сторона & Потреба & Як симулятор її задовольняє \\
    \midrule
    Студенти спеціальності 122 «Комп'ютерні науки» & Зрозуміти внутрішню роботу алгоритму ШПФ, побачити паралелізм і токенну динаміку & Анімація руху токенів по DFG, крок-за-кроком режим, теплова карта активності вузлів \\[2pt]
    Викладачі курсів ЦОС та паралельних обчислень & Демонстраційний інструмент для лекційних і лабораторних занять & Веб-застосунок без встановлення, відтворюваність на будь-якій ОС \\[2pt]
    Дослідники dataflow-архітектур & Перевірка гіпотез щодо планування та критичного шляху графа ШПФ & Метрики throughput/latency/queue, event log, налаштовувані затримки \\[2pt]
    Інженери DSP/FPGA на ранній стадії проєктування & Попередня оцінка потокової моделі перед апаратною реалізацією & Параметричне налаштування \texttt{latency} і \texttt{delay}, експорт подій \\
    \bottomrule
  \end{tabular}
\end{table}

\subsubsection*{Основні цілі розробки}

На основі проаналізованої предметної області та потреб
зацікавлених сторін визначено ієрархію цілей.

\textbf{Генеральна ціль:} створити програмний симулятор, що
відтворює потокове виконання 16-точкового алгоритму ШПФ
radix-4 $4\times4$ та забезпечує його інтерактивну графічну
візуалізацію у вигляді анімованого DFG зі спостереженням токенів,
метрик і наскрізної затримки.

Підцілі, що випливають із генеральної:
\begin{enumerate}
  \item формалізувати граф потоку даних для алгоритму radix-4
        $4\times4$ з вузлами типів \texttt{source}, \texttt{dft4},
        \texttt{twiddle}, \texttt{sink} та дугами із параметричними
        затримками;
  \item реалізувати ядро симуляції з правилом активації токенів,
        FIFO-чергами та керованим логічним часом;
  \item забезпечити інтерактивну візуалізацію руху токенів,
        теплову карту активності вузлів і спектральний вихід;
  \item виконати числову верифікацію результатів шляхом порівняння
        з еталонним обчисленням ДПФ.
\end{enumerate}

\subsubsection*{Обмеження}

Розробка виконується у таких обмеженнях:

\begin{itemize}
  \item \textit{технічні} --- цільова платформа виконання: сучасний
        веб-браузер (Chrome/Firefox/Edge) з підтримкою JavaScript
        ES2020; відсутність серверної частини (повністю клієнтська
        симуляція);
  \item \textit{ресурсні} --- розробка силами одного виконавця
        у межах бакалаврської кваліфікаційної роботи;
  \item \textit{часові} --- термін виконання згідно з календарним
        графіком кафедри;
  \item \textit{алгоритмічні} --- фіксований розмір перетворення
        $N = 16$ як базовий демонстраційний розмір; розширення на
        довільні $N = 4^t$ залишається напрямком подальшого розвитку;
  \item \textit{правові} --- використання лише програмних бібліотек
        з відкритими ліцензіями (MIT, Apache~2.0, BSD) без
        комерційних компонентів.
\end{itemize}

\subsubsection*{Критерії успіху}

Результат роботи вважається досягнутим, якщо одночасно виконуються
такі вимірювані критерії:

\begin{itemize}
  \item максимальна відхилення вихідного спектра симулятора від
        еталонного ДПФ за модулем не перевищує $\varepsilon = 10^{-9}$;
  \item чотири вузли \texttt{dft4} першого ярусу виконуються
        паралельно, що підтверджується одночасними подіями
        \texttt{onFire} у журналі симуляції;
  \item користувач може керувати швидкістю симуляції у діапазоні
        не менш ніж $0.1\times$--$10\times$ без втрати коректності;
  \item час відгуку інтерфейсу на користувацькі дії (пуск, пауза,
        зміна швидкості) не перевищує 100~мс.
\end{itemize}


% ─────────────────────────────────────────────────
\subsection{Етичні, правові та регуляторні аспекти}

\subsubsection*{Характер даних, що опрацьовуються}

Розроблюваний симулятор працює з синтетичними тестовими сигналами
(синусоїдами, імпульсами, шумом), що генеруються безпосередньо у
браузері користувача. Проєкт \emph{не опрацьовує} персональних,
медичних, біометричних або фінансових даних, не передає жодних
даних на зовнішні сервери та не використовує моделей машинного
навчання для ухвалення рішень. Унаслідок цього більшість вимог
Загального регламенту захисту даних ЄС (GDPR) та Закону України
«Про захист персональних даних» (№~2297-VI) не поширюються на
продукт у прямій формі. Водночас базові принципи мінімізації даних
і прозорості обробки враховані: застосунок не веде телеметрії,
не встановлює ідентифікаційних cookies та працює повністю
на стороні клієнта.

\subsubsection*{Академічна доброчесність}

Реалізація симулятора здійснюється з дотриманням політики
академічної доброчесності Національного університету «Львівська
політехніка». Усі використані зовнішні джерела (публікації,
бібліотеки, референтні реалізації) мають явні посилання у списку
використаних джерел. Обчислювальне ядро симулятора написане
самостійно; еталонна бібліотека \texttt{fft.js} використовується
виключно як референтна для верифікації (підрозд.~4.2) і не
залучається до шляху виконання основного алгоритму. Плагіат,
фальсифікація та фабрикація результатів не допускаються.

\subsubsection*{Ліцензійні та правові вимоги до програмного забезпечення}

Робота базується на програмних компонентах із відкритими ліцензіями:
Next.js (MIT), React (MIT), React Flow (MIT), Zustand (MIT),
TypeScript (Apache~2.0). Такі ліцензії допускають некомерційне
навчальне використання, модифікацію та розповсюдження вихідного
коду за умови збереження оригінальних повідомлень про авторство.
Розроблений у цій роботі вихідний код оприлюднюється під відкритою
ліцензією, що забезпечує відтворюваність результатів та можливість
їх повторного використання у подальших дослідженнях і навчальних
курсах.

\subsubsection*{Відповідність галузевим стандартам}

При розробці враховуються такі чинні стандарти та технічні
регламенти:

\begin{itemize}
  \item \textbf{ISO/IEC 25010:2011} --- модель якості програмного
        забезпечення; для симулятора релевантними є характеристики
        \textit{функціональна придатність}, \textit{продуктивність}
        (час відгуку $<100$~мс), \textit{зручність використання}
        та \textit{надійність} (детермінованість SDF-графа);
  \item \textbf{IEEE 754} --- стандарт подання чисел з плаваючою
        крапкою; визначає точність обчислення комплексних операцій
        у подвійній точності (double, $\varepsilon_{\mathrm{mach}}
        \approx 2.22 \cdot 10^{-16}$);
  \item \textbf{ECMAScript 2020 (ES11)} --- стандартизація мови
        JavaScript, що забезпечує переносимість коду між сучасними
        браузерами;
  \item \textbf{ДСТУ 3008:2015} --- вимоги до оформлення звітів у
        сфері науки і техніки, що застосовуються до пояснювальної
        записки БКР;
  \item \textbf{ДСТУ ГОСТ 7.1:2006} --- правила бібліографічного
        опису, що застосовуються у списку використаних джерел.
\end{itemize}

Оскільки проєкт не використовує персональних даних і не залучає
моделей ШІ для ухвалення рішень, питання упередженості моделей
(bias/fairness), галюцинацій LLM та конфіденційності запитів у
межах цієї роботи є нерелевантними.


% ─────────────────────────────────────────────────
\subsection{Постановка задачі та формулювання вимог}

\subsubsection*{Формальна постановка задачі}

Нехай задано комплекснозначну вхідну послідовність
$\{x_n\}_{n=0}^{N-1}$ довжиною $N = 16$, подану у подвійній
точності формату IEEE~754. Потрібно розробити програмний
симулятор, який:

\begin{enumerate}
  \item обчислює дискретне перетворення Фур'є
        \begin{equation*}
          X_k = \sum_{n=0}^{N-1} x_n \, e^{-j 2\pi n k / N},
          \quad k = 0, 1, \ldots, N-1,
        \end{equation*}
        за алгоритмом radix-4 з двовимірною $4 \times 4$
        декомпозицією, реалізованим у моделі потоку даних;
  \item формалізує процес обчислення у вигляді орієнтованого
        графа потоку даних $G = (V, E)$, де $V$ --- множина вузлів
        типів \{\texttt{source}, \texttt{dft4}, \texttt{twiddle},
        \texttt{sink}\}, $E$ --- множина спрямованих дуг із
        параметричними затримками;
  \item виконує симуляцію графа з дотриманням правила активації
        (вузол спрацьовує тоді і лише тоді, коли на всіх його
        вхідних каналах є токени) та з урахуванням індивідуальних
        затримок \texttt{latency} вузлів і \texttt{delay} дуг;
  \item забезпечує інтерактивну візуалізацію процесу обчислення
        та подає кінцевий спектр $\{X_k\}$ у вигляді амплітудної
        діаграми.
\end{enumerate}

\subsubsection*{Вхідні та вихідні дані}

\textbf{Вхідні дані симулятора:}
\begin{itemize}
  \item послідовність $\{x_n\}_{n=0}^{15}$ --- 16~комплексних
        відліків сигналу, задається користувачем або генерується
        вбудованим джерелом (синусоїда, одиничний імпульс,
        білий шум);
  \item параметри симуляції --- затримки вузлів
        ($L_{\mathrm{dft4}}$, $L_{\mathrm{twiddle}}$), затримка
        дуг ($\Delta_{\mathrm{edge}}$), множник швидкості
        (speed multiplier);
  \item команди керування --- пуск, пауза, покроковий режим,
        скидання стану.
\end{itemize}

\textbf{Вихідні дані симулятора:}
\begin{itemize}
  \item спектр $\{X_k\}_{k=0}^{15}$ --- 16~комплексних відліків,
        що візуалізуються як амплітуда $|X_k|$ і фаза $\arg X_k$;
  \item журнал подій (event log) --- часові мітки подій
        \texttt{onToken}, \texttt{onFire}, \texttt{onOutput} для
        кожного вузла;
  \item метрики роботи графа --- пропускна здатність (throughput),
        наскрізна затримка (latency), розмір черг (queue size),
        теплова карта активності вузлів;
  \item звіт про верифікацію --- максимальне відхилення $\max_k
        |X_k^{\mathrm{sim}} - X_k^{\mathrm{ref}}|$ відносно
        еталонного ДПФ.
\end{itemize}

\subsubsection*{Функціональні вимоги}

Функціональні вимоги до симулятора зведено у табл.~\ref{tab:func_req}.

\begin{table}[H]
  \caption{Функціональні вимоги до симулятора}
  \label{tab:func_req}
  \centering
  \small
  \begin{tabular}{p{1.4cm}p{4.5cm}p{7.5cm}}
    \toprule
    ID & Назва & Опис \\
    \midrule
    F1 & Побудова DFG & Автоматична генерація графа потоку даних для ШПФ radix-4 $4\times4$ при $N = 16$ \\
    F2 & Симуляція виконання & Крокування часу з доставкою токенів, перевіркою правила активації та виконанням вузлів \\
    F3 & Введення сигналу & Вибір типу вхідного сигналу (синусоїда, імпульс, шум) та його параметрів \\
    F4 & Керування симуляцією & Пуск, пауза, покроковий режим, скидання, зміна множника швидкості \\
    F5 & Візуалізація токенів & Анімоване переміщення токенів по дугах з відображенням значень \\
    F6 & Відображення метрик & Показ throughput, latency, queue size, теплової карти активності \\
    F7 & Спектральний вихід & Відображення амплітуди та фази $\{X_k\}$ у вигляді гістограми \\
    F8 & Верифікація результату & Порівняння з еталонним ДПФ і виведення максимального відхилення \\
    F9 & Налаштування параметрів & Редагування затримок \texttt{latency}, \texttt{delay} без перезавантаження сторінки \\
    \bottomrule
  \end{tabular}
\end{table}

\subsubsection*{Нефункціональні вимоги}

Нефункціональні вимоги регламентують якісні характеристики системи
(табл.~\ref{tab:nonfunc_req}) у відповідності до ISO/IEC~25010.

\begin{table}[H]
  \caption{Нефункціональні вимоги до симулятора}
  \label{tab:nonfunc_req}
  \centering
  \small
  \begin{tabular}{p{1.4cm}p{3.8cm}p{8.2cm}}
    \toprule
    ID & Характеристика & Вимога \\
    \midrule
    NF1 & Точність обчислень & $\max_k |X_k^{\mathrm{sim}} - X_k^{\mathrm{ref}}| \le 10^{-9}$ \\
    NF2 & Час відгуку UI & Реакція на користувацьку дію $\le 100$~мс \\
    NF3 & Частота кадрів & Стабільні 60~fps при швидкості симуляції $\le 5\times$ \\
    NF4 & Портативність & Запуск у Chrome~$\ge$~90, Firefox~$\ge$~88, Edge~$\ge$~90 \\
    NF5 & Детермінованість & Однаковий вхід і параметри дають ідентичний журнал подій \\
    NF6 & Відкритий код & Розповсюдження під ліцензією MIT або Apache~2.0 \\
    NF7 & Без встановлення & Робота у браузері без серверної частини та БД \\
    NF8 & Зручність & Основні операції доступні не глибше двох кліків від головного екрану \\
    \bottomrule
  \end{tabular}
\end{table}

\subsubsection*{Діаграма прецедентів (Use Case)}

У системі виділяються два актори --- \emph{Студент} (первинний
користувач, спостерігач) і \emph{Викладач} (вторинний користувач,
може налаштовувати сценарії демонстрації). Основні прецеденти та
їхні зв'язки наведено у табл.~\ref{tab:usecases}.

\begin{table}[H]
  \caption{Перелік прецедентів (Use Cases) симулятора}
  \label{tab:usecases}
  \centering
  \small
  \begin{tabular}{p{1.2cm}p{4.0cm}p{2.0cm}p{6.5cm}}
    \toprule
    ID & Прецедент & Актор & Стислий опис \\
    \midrule
    UC1 & Завантажити готовий сценарій & Студент, Викладач & Обрати один із передустановлених входів (синусоїда, імпульс, шум) \\
    UC2 & Задати власний вхідний сигнал & Студент & Ввести 16~комплексних відліків вручну \\
    UC3 & Запустити симуляцію & Студент & Ініціювати виконання DFG з пуску до завершення \\
    UC4 & Керувати темпом симуляції & Студент & Пауза, step-by-step, зміна множника швидкості \\
    UC5 & Спостерігати рух токенів & Студент & Переглядати анімацію токенів по дугах у реальному часі \\
    UC6 & Переглянути спектр результату & Студент & Побачити $|X_k|$ і $\arg X_k$ у вигляді діаграми \\
    UC7 & Порівняти з еталоном & Студент, Викладач & Запустити \texttt{compareWithFFT()} і побачити похибку \\
    UC8 & Змінити параметри затримок & Викладач & Налаштувати \texttt{latency} вузлів і \texttt{delay} дуг \\
    UC9 & Експортувати журнал подій & Викладач & Зберегти event log у файл для аналізу \\
    \bottomrule
  \end{tabular}
\end{table}

\emph{Підпорядкування прецедентів:} UC3 містить (\textit{include})
UC5 (рух токенів автоматично супроводжує виконання); UC7
розширює (\textit{extend}) UC3 та може бути викликаний після
завершення симуляції; UC1 і UC2 є взаємовиключними альтернативами
підготовки вхідних даних.

Описана постановка задачі та перелік вимог утворюють технічне
завдання, яке далі конкретизується у розділі~2 через вибір методів
і технологій реалізації, проєктування архітектури та структур
даних.


% ─────────────────────────────────────────────────
\subsection*{Висновки до розділу~1}

У першому розділі виконано аналіз предметної області, огляд сучасних
підходів, системний аналіз задачі, розгляд етичних та правових
аспектів і формалізовано постановку задачі.

Зафіксовано \emph{стан предметної області:} алгоритм ШПФ radix-4
$4\times4$ для $N = 16$ має достатньо розроблений математичний
апарат і формальні потокові моделі (DFG, SDF), однак у практичному
інструментарії ці рівні описані відокремлено. Виробничі бібліотеки
(FFTW, ARM~CMSIS-DSP, NVIDIA~CuFFT, SPIRAL, Intel~IPP) оптимізовані
виключно для продуктивності та не надають засобів інтроспекції.
Методи візуального моделювання (UML, IDEF0, DFG/SDF, сигнальні
графи) залишаються переважно теоретичними та не реалізуються як
інтерактивні веб-симулятори з числовою верифікацією.

Виявлено \emph{прогалину:} жодне з розглянутих рішень не поєднує
одночасно достовірного обчислення ШПФ за формальною SDF-моделлю,
інтерактивної динамічної візуалізації руху токенів і веб-доступності
без комерційної ліцензії. Саме ця незаповнена ніша визначає предмет
цієї роботи.

Сформульовано \emph{задачу} розробки програмного симулятора потокового
обчислювача 16-точкового ШПФ з анімованим DFG, визначено вхідні та
вихідні дані, технічні, ресурсні й правові обмеження та критерії
успіху. Розроблено перелік із 9~функціональних (F1--F9) і 8~нефункціональних
(NF1--NF8) вимог, побудовано перелік із 9~прецедентів (UC1--UC9),
що охоплюють основні сценарії використання симулятора студентами
та викладачами. Ключові вимірювані критерії --- точність $\varepsilon
\le 10^{-9}$, час відгуку $\le 100$~мс, стабільні 60~fps при
швидкості до $5\times$ --- підлягають експериментальній перевірці
у розділі~4.

Розглянуто етичні та правові аспекти: робота не опрацьовує
персональних даних, базується на відкритих ліцензіях (MIT,
Apache~2.0) і відповідає стандартам ISO/IEC~25010, IEEE~754,
ECMAScript~2020, ДСТУ~3008:2015 та ДСТУ~ГОСТ~7.1:2006.

Математичне та алгоритмічне обґрунтування (виведення формул
radix-4 $4\times4$, матричні факторизації, формалізація DFG/SDF)
виноситься до розділу~2, де воно подається у складі проєктного
обґрунтування обраних методів і архітектури системи.
